\chapter{Anhang}

% NEU (21.10.): aus Abschnitt 3.4.1 verschoben (Entscheidung E3), damit die Gewichtung im Ansatzkapitel stimmt; Text unverändert bis auf den Einleitungssatz
\section{Gegenüberstellung der Ähnlichkeitsmaße}
\label{Anhang_Aehnlichkeit}

Die Wahl des Ähnlichkeitsmaßes in Abschnitt~\ref{subsec:ta-erschliessen} stützt sich auf die folgende Abwägung der in Abschnitt~\ref{subsec:embeddings} eingeführten Maße. Ein Skalarprodukt ohne vorherige Normalisierung ist nicht rangäquivalent zur Kosinusähnlichkeit, da der Zusammenhang zur Distanz nur bei konstanter Vektorlänge besteht und die Längen realer Repräsentationen erheblich schwanken \cite{shrivastava_asymmetric_2014}. Es bevorzugt folglich Dokumente mit großer Vektornorm unabhängig von deren inhaltlicher Passung. Die Manhattan-Distanz bleibt auch nach Normalisierung rangverschieden, da sie sich nicht als Funktion des Skalarprodukts allein darstellen lässt, und ist für hochdimensionale dichte Repräsentationen nicht etabliert. Der Jaccard-Koeffizient misst die Überschneidung zweier Mengen \cite{manning_introduction_2008} und setzt damit diskrete Merkmale voraus. Angewandt auf Bezeichner würde er auf die Übereinstimmung von Zeichenketten zurückfallen und damit das Vokabularproblem reproduzieren, dessen Umgehung den Einsatz von Embeddings begründet. Gleiches gilt für Maße der Zeichenkettenähnlichkeit wie die Editierdistanz \cite{manning_introduction_2008}, welche die Anzahl der zur Überführung zweier Zeichenketten erforderlichen Operationen zählt. Sie eignen sich zur Erkennung von Schreibvarianten innerhalb eines bekannten Bezeichnervorrats, nicht jedoch zur Überbrückung der Lücke zwischen einer fachlichen Anfrage und einem technischen Kürzel. Tabelle~\ref{tab:aehnlichkeitsmasse} stellt die betrachteten Maße gegenüber.

\begin{table}[htbp]
\centering
\caption{Gegenüberstellung der betrachteten Ähnlichkeits- und Distanzmaße}
\label{tab:aehnlichkeitsmasse}
\small
\begin{tabular}{p{3.0cm}p{3.4cm}p{6.4cm}}
\hline
\textbf{Maß} & \textbf{Gegenstand} & \textbf{Einordnung für diese Arbeit} \\
\hline
Kosinusähnlichkeit & Winkel zwischen dichten Vektoren & Verwendet; längeninvariant und im Information Retrieval etabliert \\
Skalarprodukt (normalisiert) & Winkel, implizit & Verwendet; auf Einheitsvektoren rangäquivalent zur Kosinusähnlichkeit und effizienter auswertbar \\
Skalarprodukt (unnormalisiert) & Winkel und Länge & Verworfen; bevorzugt Vektoren großer Norm unabhängig von der inhaltlichen Passung \\
Euklidische Distanz & Abstand im Raum & Nicht verwendet; auf Einheitsvektoren rangäquivalent, jedoch rechenaufwendiger \\
Manhattan-Distanz & Abstand entlang der Achsen & Verworfen; auch normalisiert nicht rangäquivalent, für dichte Repräsentationen unüblich \\
Jaccard-Koeffizient & Mengenüberschneidung & Verworfen; setzt diskrete Merkmale voraus und fällt auf lexikalische Übereinstimmung zurück \\
Editierdistanz & Zeichenketten\-abstand & Verworfen; erkennt Schreibvarianten, überbrückt jedoch keine Bedeutungsunterschiede \\
\hline
\end{tabular}
\end{table}

\section{Ergänzende Diagramme}
\label{Anhang_Diagramme}

Hier könnte noch was stehen. Z.B. zusätzliche Diagramme, Gleichungen
oder erklärender Text, die für den Lesefluss im Haupttext hinterlich
sind, aber interessante zusätzliche Einblicke liefern.

\section{Ergänzende Codeschnipsel}
\label{Anhang_Code}

Hier könnten auch Code-Snippets stehen, die kein Mensch liest, die
aber ungeheueren Eindruck machen.
